\chapter[18]{Geometric Tools and Point Picking Methods}
\label{ch:chow-iii-geometric-tools-point-picking}

Write \(d_g\), \(B_g(x,r)\), and \(L_g\) for distance, balls, and length.
For a time-dependent metric, \(\mathcal Z(t_0)\) is the compact set of
unit-speed \(g(t_0)\)-minimizing geodesics from \(x_0\) to \(x\).

\section{Estimates for changing distances}

\begin{lemma}[Time derivative of the distance function]
\label{lem:chow-iii-forward-distance-derivative}
\source{chow-et-al-ricci-flow-part-iii:page-0062,chow-et-al-ricci-flow-part-iii:page-0063}
\dcref{ch18:18.1}
\group{Changing Distances}

Let \(g(t)\), \(t\in[0,T)\), be a smooth family of complete metrics.  In the
sense of the lower forward derivative,
\[
 \left.\frac{\partial^-}{\partial t}\right|_{t=t_0}d_{g(t)}(x,x_0)
 =\min_{\eta\in\mathcal Z(t_0)}\frac12\int_\eta
 \partial_tg(t_0)(\eta',\eta')\,\dd s .
\]
\end{lemma}
\begin{proof}
\sketch For a fixed \(\eta\in\mathcal Z(t_0)\), comparison with its length at nearby
times and the first variation formula give the upper bound.  Conversely,
choose minimizing geodesics \(\eta_i\) at times \(t_0+h_i\), where
\(h_i\downarrow0\).  Compactness of \(\mathcal Z(t_0)\) gives a subsequence
converging to \(\eta_\infty\in\mathcal Z(t_0)\).  Applying the mean value
theorem to \(L_{g(t)}(\eta_i)\) and passing to the limit gives the reverse
inequality, hence equality.
\end{proof}

\begin{corollary}[Time derivative of distance under Ricci flow]
\label{cor:chow-iii-forward-distance-derivative-ricci-flow}
\source{chow-et-al-ricci-flow-part-iii:page-0063}
\uses{lem:chow-iii-forward-distance-derivative}
\dcref{ch18:18.2}
\group{Changing Distances}


If \(g(t)\) is a complete Ricci flow, then
\[
 \left.\frac{\partial^-}{\partial t}\right|_{t=t_0}d_{g(t)}(x,x_0)
 =\min_{\eta\in\mathcal Z(t_0)}
 \left(-\int_\eta\Ric_{g(t_0)}(\eta',\eta')\,\dd s\right).
\]
\end{corollary}
\begin{proof}
\sketch Substitute \(\partial_tg=-2\Ric\) in
Lemma~\ref{lem:chow-iii-forward-distance-derivative}.
\end{proof}

\begin{lemma}[Backward lower derivative of the distance function]
\label{lem:chow-iii-backward-distance-derivative}
\source{chow-et-al-ricci-flow-part-iii:page-0063,chow-et-al-ricci-flow-part-iii:page-0064}
\dcref{ch18:18.4}
\group{Changing Distances}

If \(g(t)\) is a smooth family of complete metrics, then, in the sense of the
lower backward derivative,
\[
 \left.\frac{\partial_-}{\partial t}\right|_{t=t_0}d_{g(t)}(x,x_0)
 \ge\frac12\max_{\eta\in\mathcal Z(t_0)}
 \int_\eta\partial_tg(t_0)(\eta',\eta')\,\dd s.
\]
\end{lemma}
\begin{proof}
\sketch For every \(g(t_0)\)-minimizing geodesic \(\eta\), one has
\(L_{g(t)}(\eta)\ge d_{g(t)}(x,x_0)\) for earlier times and equality at
\(t_0\).  Taking the lower backward derivative and applying the first
variation formula for length gives the asserted inequality for each \(\eta\);
maximizing over the compact set \(\mathcal Z(t_0)\) completes the proof.
\end{proof}

\begin{definition}[Laplacian upper bound in the barrier sense]
\label{def:chow-iii-laplacian-upper-barrier}
\source{chow-et-al-ricci-flow-part-iii:page-0064}
\dcref{ch18:18.5}
\group{Distance Barriers}

Let \(\varphi\) be continuous near \(x\).  The inequality
\(\Delta\varphi(x)\le C\) holds in the barrier sense if, for every
\(\varepsilon>0\), there is a \(C^2\) function \(\psi\) near \(x\) such that
\(\psi(x)=\varphi(x)\), \(\psi\ge\varphi\) nearby, and
\(\Delta\psi(x)\le C+\varepsilon\).
\end{definition}

\begin{lemma}[Upper bound for the Laplacian of distance]
\label{lem:chow-iii-distance-laplacian-upper-bound}
\source{chow-et-al-ricci-flow-part-iii:page-0064,chow-et-al-ricci-flow-part-iii:page-0065,chow-et-al-ricci-flow-part-iii:page-0066}
\uses{def:chow-iii-laplacian-upper-barrier}
\dcref{ch18:18.6}
\group{Distance Barriers}


Let \((\mathcal M^n,g)\) be complete, let \(x\ne x_0\), and let
\(\gamma:[0,s_0]\to\mathcal M\) be a unit-speed minimizing geodesic from
\(x_0\) to \(x\).  For every continuous piecewise smooth
\(\zeta:[0,s_0]\to[0,1]\) with \(\zeta(0)=0\) and \(\zeta(s_0)=1\),
\[
 \Delta_gd_g(x_0,x)\le\int_0^{s_0}
 \bigl((n-1)(\zeta')^2-\zeta^2\Ric(\gamma',\gamma')\bigr)\,\dd s
\]
in the barrier sense.
\end{lemma}
\begin{proof}
\sketch Vary \(\gamma\) through paths whose endpoint moves by the exponential map at
\(x\), using variation fields \(\zeta E_i\), where the \(E_i\) are parallel
orthonormal fields normal to \(\gamma'\).  Their lengths define a smooth upper
barrier for distance.  Summing the second variation formula over the
\(n-1\) normal directions gives
\[
 \sum_{i=1}^{n-1}\int_0^{s_0}
 \bigl((\zeta')^2-\zeta^2\langle R(\gamma',E_i)E_i,\gamma'\rangle\bigr)\,\dd s,
\]
which is the displayed expression after taking the trace of curvature.
\end{proof}

\begin{theorem}[Estimate for changing distances]
\label{thm:chow-iii-changing-distances}
\source{chow-et-al-ricci-flow-part-iii:page-0066,chow-et-al-ricci-flow-part-iii:page-0067,chow-et-al-ricci-flow-part-iii:page-0068,chow-et-al-ricci-flow-part-iii:page-0069}
\uses{cor:chow-iii-forward-distance-derivative-ricci-flow, lem:chow-iii-backward-distance-derivative, lem:chow-iii-distance-laplacian-upper-bound}
\dcref{ch18:18.7}
\group{Changing Distances}


Let \((\mathcal M^n,g(t))\), \(t\in[0,T)\), be a complete Ricci flow.
If \(\Ric(\cdot,t_0)\le(n-1)K\) on \(B_{g(t_0)}(x_0,r_0)\), then outside
that ball
\[
 (\partial_t-\Delta)d_{g(t)}(x,x_0)|_{t=t_0}
 \ge-(n-1)\left(\frac23Kr_0+\frac1{r_0}\right)
\]
in the barrier sense.  If the same Ricci bound holds on the union of the
\(r_0\)-balls about \(x_0,x_1\), then
\[
 \partial_td_{g(t)}(x_0,x_1)|_{t=t_0}\ge
 \begin{cases}
 -2(n-1)(\frac23Kr_0+r_0^{-1}),&d_{g(t_0)}(x_0,x_1)\ge2r_0,\\
 -2(n-1)Kr_0,&d_{g(t_0)}(x_0,x_1)<2r_0.
 \end{cases}
\]
\end{theorem}
\begin{proof}
\sketch Choose a minimizing geodesic realizing the appropriate one-sided distance
derivative.  In Lemma~\ref{lem:chow-iii-distance-laplacian-upper-bound}, take
\(\zeta(s)=s/r_0\) on \([0,r_0]\) and \(1\) thereafter.  Integrating
\((n-1)K(1-s^2/r_0^2)+(n-1)r_0^{-2}\) gives the first estimate.  For two
points at distance at least \(2r_0\), apply this argument from both endpoints
to an interior point of a minimizing geodesic and use that the sum of the two
distance functions has a minimum there.  If their distance is less than
\(2r_0\), integrate the Ricci upper bound directly along the minimizing
geodesic.  These are precisely the two asserted cases.
\end{proof}

\begin{proposition}[Estimate for Ricci curvature along a stable geodesic]
\label{prop:chow-iii-ricci-integral-stable-geodesic}
\source{chow-et-al-ricci-flow-part-iii:page-0069}
\uses{lem:chow-iii-distance-laplacian-upper-bound}
\dcref{ch18:18.8}
\group{Changing Distances}


Let \(\gamma:[0,L]\to\mathcal M\) be a stable unit-speed geodesic with
\(L\ge2r\), and suppose
\(\Ric\le(n-1)K\) on \(B(\gamma(0),r)\cup B(\gamma(L),r)\).  Then
\[
 \int_\gamma\Ric(\gamma',\gamma')\,\dd s
 \le2(n-1)\left(\frac23Kr+\frac1r\right).
\]
In particular, for \(r=K^{-1/2}\) this is at most
\(\frac{10}{3}(n-1)\sqrt K\).
\end{proposition}
\begin{proof}
\sketch Use in the index form the piecewise linear test function which rises from
zero to one over the initial segment of length \(r\), stays one, and falls to
zero over the final segment of length \(r\).  Stability says the sum of the
normal index forms is nonnegative.  Rearranging and integrating the endpoint
Ricci bounds gives the displayed estimate; substitution of \(r=K^{-1/2}\)
gives the final constant.
\end{proof}

\begin{remark}[Changing-distance bounds for ancient solutions]
\label{rem:chow-iii-changing-distances-ancient-solutions}
\source{chow-et-al-ricci-flow-part-iii:page-0070}
\uses{thm:chow-iii-changing-distances, prop:chow-iii-ricci-integral-stable-geodesic}
\dcref{ch18:18.9}
\group{Changing Distances}


The preceding bounds on the time derivative of distance are particularly
useful in the study of ancient Ricci-flow solutions.
\end{remark}

\section{Spatial point picking methods}

\begin{theorem}[Point picking when the asymptotic scalar curvature ratio is infinite]
\label{thm:chow-iii-point-picking-infinite-ascr}
\source{chow-et-al-ricci-flow-part-iii:page-0070,chow-et-al-ricci-flow-part-iii:page-0071}
\dcref{ch18:18.10}
\group{Spatial Point Picking}

Let \((\mathcal M^n,g,O)\) be complete and noncompact with bounded scalar
curvature and \(\operatorname{ASCR}(g)=\infty\).  There are points
\(x_i\to\infty\), numbers \(\varepsilon_i\to0\), and radii \(r_i>0\) such
that the balls \(B(x_i,r_i)\) are disjoint and
\[
 \sup_{B(x_i,r_i)}R\le(1+\varepsilon_i)R(x_i),\qquad
 R(x_i)r_i^2\to\infty,\qquad \frac{d(x_i,O)}{r_i}\to\infty.
\]
\end{theorem}

\begin{remark}[Rescaled metrics have bounded scalar curvature on large balls]
\label{rem:chow-iii-rescaled-point-picked-balls}
\source{chow-et-al-ricci-flow-part-iii:page-0071}
\uses{thm:chow-iii-point-picking-infinite-ascr}
\dcref{ch18:18.11}
\group{Spatial Point Picking}


For \(g_i=R(x_i)g\), one has \(R_{g_i}(x_i)=1\) and
\[
 \sup_{B_{g_i}(x_i,\widetilde r_i)}R_{g_i}\le1+\varepsilon_i,
 \qquad \widetilde r_i=R(x_i)^{1/2}r_i\to\infty,
\]
while \(d_{g_i}(x_i,O)\to\infty\).
\end{remark}

\begin{proposition}[Unbounded-curvature point picking, version I]
\label{prop:chow-iii-unbounded-curvature-point-picking-one}
\source{chow-et-al-ricci-flow-part-iii:page-0071,chow-et-al-ricci-flow-part-iii:page-0072}
\dcref{ch18:18.12}
\group{Spatial Point Picking}

If \((\mathcal M^n,g,O)\) is complete and noncompact with
\(\sup_{\mathcal M} R=\infty\), there are disjoint balls \(B(x_i,r_i)\), with
\(r_i\in(0,1]\) and \(d(x_i,O)\to\infty\), such that
\[
 \sup_{B(x_i,r_i)}R\le4R(x_i),\qquad R(x_i)r_i^2\to\infty.
\]
\end{proposition}
\begin{proof}
\sketch Choose \(y_i\) with \(R(y_i)\to\infty\), and maximize
\(R(x)d(x,\partial B(y_i,2))^2\) on \(B(y_i,2)\).  If \(x_i\) is a maximizer,
set \(r_i=\frac12d(x_i,\partial B(y_i,2))\).  Maximality gives the factor-four
bound and \(R(x_i)r_i^2\ge R(y_i)\).  Properness of scalar curvature along the
chosen sequence gives \(d(x_i,O)\to\infty\), and a subsequence has disjoint
balls because \(r_i\le1\).
\end{proof}

\begin{proposition}[Unbounded-curvature point picking, version II]
\label{prop:chow-iii-unbounded-curvature-point-picking-two}
\source{chow-et-al-ricci-flow-part-iii:page-0072,chow-et-al-ricci-flow-part-iii:page-0073}
\dcref{ch18:18.13}
\group{Spatial Point Picking}

Let \(\varepsilon_i>0\) be bounded.  Under the hypotheses of
Proposition~\ref{prop:chow-iii-unbounded-curvature-point-picking-one}, there
are disjoint balls \(B(x_i,r_i)\), with \(r_i\in(0,1]\) and
\(d(x_i,O)\to\infty\), such that
\[
 \sup_{B(x_i,r_i)}R\le(1+\varepsilon_i)R(x_i),
 \qquad R(x_i)r_i^2\to\infty.
\]
\end{proposition}
\begin{proof}
\sketch Choose \(y_i\) so that \(R(y_i)\varepsilon_i^2\to\infty\), maximize
\(R(x)d(x,\partial B(y_i,1))^2\), and write the maximizing boundary distance
as \(\widetilde r_i\).  Set
\[
 \widetilde\varepsilon_i=
 \frac{\varepsilon_i}{\sqrt{1+\varepsilon_i}(1+\sqrt{1+\varepsilon_i})},
 \qquad r_i=\widetilde\varepsilon_i\widetilde r_i.
\]
Maximality then yields \(R\le(1+\varepsilon_i)R(x_i)\) on \(B(x_i,r_i)\),
and \(R(x_i)r_i^2\to\infty\).  The remaining conclusions follow as in
Proposition~\ref{prop:chow-iii-unbounded-curvature-point-picking-one}.
\end{proof}

\begin{remark}[Rescaling the improved point-picked sequence]
\label{rem:chow-iii-improved-point-picking-rescaling}
\source{chow-et-al-ricci-flow-part-iii:page-0072}
\uses{prop:chow-iii-unbounded-curvature-point-picking-two, rem:chow-iii-rescaled-point-picked-balls}
\dcref{ch18:18.14}
\group{Spatial Point Picking}


The rescaling conclusion of
Remark~\ref{rem:chow-iii-rescaled-point-picked-balls} applies equally to the
sequence in Proposition~\ref{prop:chow-iii-unbounded-curvature-point-picking-two}.
\end{remark}

\begin{lemma}[Point picking when scalar-curvature change is unbounded]
\label{lem:chow-iii-point-picking-unbounded-change}
\source{chow-et-al-ricci-flow-part-iii:page-0073,chow-et-al-ricci-flow-part-iii:page-0074}
\dcref{ch18:18.15}
\group{Spatial Point Picking}

Let \((\mathcal N_k^n,h_k,O_k)\) be complete pointed manifolds with
\(R_{h_k}(O_k)=1\).  If, for some \(D>0\),
\(\sup_{B_{h_k}(O_k,D)}R_{h_k}\to\infty\), then there are
\(w_k\in B_{h_k}(O_k,D+1)\) and \(s_k\in(0,D+1)\) such that
\[
 R_{h_k}(w_k)s_k^2\to\infty,
 \qquad \sup_{B_{h_k}(w_k,s_k)}R_{h_k}\le2R_{h_k}(w_k).
\]
\end{lemma}
\begin{proof}
\sketch Maximize
\(F_k(w)=R_{h_k}(w)(D+1-d_{h_k}(w,O_k))^2\) on the compact closed ball.
At a maximizer \(w_k\), set
\(s_k=(1-2^{-1/2})(D+1-d(w_k,O_k))\).  The hypothesis forces
\(F_k(w_k)\to\infty\), hence the first conclusion.  For
\(w\in B(w_k,s_k)\), the triangle inequality bounds its distance to the
boundary from below by \(s_k/(\sqrt2-1)\); comparison with the maximum of
\(F_k\) gives \(R_{h_k}(w)\le2R_{h_k}(w_k)\).
\end{proof}

\begin{remark}[Curvature control on expanding rescaled balls]
\label{rem:chow-iii-unbounded-change-rescaled-balls}
\source{chow-et-al-ricci-flow-part-iii:page-0074}
\uses{lem:chow-iii-point-picking-unbounded-change}
\dcref{ch18:18.16}
\group{Spatial Point Picking}


After scaling by \(R_{h_k}(w_k)\), the balls in
Lemma~\ref{lem:chow-iii-point-picking-unbounded-change} have radii tending to
infinity and scalar curvature at most twice its value at the center.
\end{remark}

\begin{theorem}[Dimension reduction by splitting a line, version I]
\label{thm:chow-iii-dimension-reduction-one}
\source{chow-et-al-ricci-flow-part-iii:page-0075}
\dcref{ch18:18.17}
\group{Dimension Reduction}

Let \((\mathcal M^n,g(t),O)\), \(t\in(-\infty,\omega)\), be a pointed complete
noncompact ancient solution with bounded \(\Rm\ge0\).  Suppose there are
\(x_i\), \(r_i\to\infty\), and \(A_i\to\infty\) such that
\[
 d_{g(0)}(O,x_i)/r_i\ge A_i,qquad
 \Rm(\cdot,0)\le r_i^{-2}\ \text{on }B_{g(0)}(x_i,A_ir_i),
\]
and \(\operatorname{inj}_{g(0)}(x_i)\ge\iota_0r_i\).  A subsequence of
\((\mathcal M,r_i^{-2}g(r_i^2t),x_i)\) converges smoothly in the pointed sense
to a complete ancient product of a bounded nonnegatively curved
\((n-1)\)-dimensional solution and a line.
\end{theorem}

\begin{remark}[Interpretation of the first dimension-reduction criterion]
\label{rem:chow-iii-dimension-reduction-one-interpretation}
\source{chow-et-al-ricci-flow-part-iii:page-0075}
\uses{thm:chow-iii-dimension-reduction-one}
\dcref{ch18:18.18}
\group{Dimension Reduction}


The hypotheses give curvature control on large balls centered relatively far
from the origin, together with an injectivity-radius bound.  Since
\(r_i\to\infty\), the resulting limit is a blow-down limit.
\end{remark}

\begin{theorem}[Dimension reduction by splitting a line, version II]
\label{thm:chow-iii-dimension-reduction-two}
\source{chow-et-al-ricci-flow-part-iii:page-0075,chow-et-al-ricci-flow-part-iii:page-0076,chow-et-al-ricci-flow-part-iii:page-0077}
\dcref{ch18:18.19}
\group{Dimension Reduction}

Let \((\mathcal M^n,g(t))\) be a complete noncompact ancient solution with
bounded \(\Rm\ge0\), and fix \(O\in\mathcal M\).  Suppose that for some
\(C<\infty\) there are \(x_i\) and \(r_i>0\) with
\[
 d_{g(0)}(O,x_i)^2R(x_i,0)\to\infty,
 \quad r_i^2R(x_i,0)\to\infty,
 \quad \Rm(\cdot,0)\le CR(x_i,0)\ \text{on }B_{g(0)}(x_i,r_i),
\]
and \(\operatorname{inj}_{g(0)}(x_i)\ge\iota_0R(x_i,0)^{-1/2}\).  A
subsequence of the flows scaled by \(R(x_i,0)\) about \(x_i\) converges to a
complete nonflat product of a bounded nonnegatively curved
\((n-1)\)-dimensional solution and a line.
\end{theorem}
\begin{proof}
\sketch The trace Harnack inequality propagates the curvature bound backward in time,
and the injectivity-radius hypothesis gives smooth pointed compactness.  The
two divergence assumptions make the limit complete and nonflat.  Choose
successive \(x_i\) so that minimizing geodesics from \(O\) approach a ray and
the geodesics from \(x_i\) to \(x_{i+1}\) extend arbitrarily far after scaling.
Toponogov comparison and its angle monotonicity show that their two limiting
rays based at \(x_\infty\) meet at angle \(\pi\), hence form a line.  The
Cheeger--Gromoll splitting theorem, preserved along the limiting Ricci flow,
gives the asserted product.
\end{proof}

\begin{remark}[Infinite asymptotic scalar curvature ratio]
\label{rem:chow-iii-dimension-reduction-infinite-ascr}
\source{chow-et-al-ricci-flow-part-iii:page-0076}
\uses{thm:chow-iii-dimension-reduction-two}
\dcref{ch18:18.20}
\group{Dimension Reduction}


The hypotheses of Theorem~\ref{thm:chow-iii-dimension-reduction-two} imply
\(\operatorname{ASCR}(g(0))=\infty\), and its curvature bounds are comparable
to the curvature at the centers, ensuring that the limit is nonflat.
\end{remark}

\begin{corollary}[Dimension reduction for ancient kappa-solutions]
\label{cor:chow-iii-dimension-reduction-kappa-solutions}
\source{chow-et-al-ricci-flow-part-iii:page-0077}
\uses{thm:chow-iii-point-picking-infinite-ascr, thm:chow-iii-dimension-reduction-two}
\dcref{ch18:18.21}
\group{Dimension Reduction}


A noncompact ancient \(\kappa\)-solution with
\(\operatorname{ASCR}(t)=\infty\) dimension reduces to the product of a
bounded nonnegatively curved \((n-1)\)-dimensional \(\kappa\)-solution and a
line.
\end{corollary}
\begin{proof}
\sketch Theorem~\ref{thm:chow-iii-point-picking-infinite-ascr} provides the required
centers and scales.  Noncollapsing supplies the injectivity-radius estimate,
so Theorem~\ref{thm:chow-iii-dimension-reduction-two} applies.  Noncollapsing
passes to smooth pointed limits and to the product factor.
\end{proof}

\begin{remark}[Kappa-solutions with Harnack]
\label{rem:chow-iii-kappa-solution-harnack-ascr}
\source{chow-et-al-ricci-flow-part-iii:page-0077}
\dcref{ch18:18.22}
\group{Dimension Reduction}

A noncompact \(\kappa\)-solution satisfying Harnack has infinite asymptotic
scalar curvature ratio at every time.
\end{remark}

\begin{theorem}[Finitely many remote curvature bumps]
\label{thm:chow-iii-finite-remote-curvature-bumps}
\source{chow-et-al-ricci-flow-part-iii:page-0077,chow-et-al-ricci-flow-part-iii:page-0078}
\dcref{ch18:18.23}
\group{Dimension Reduction}

For every \(\varepsilon>0\) and \(n\ge2\), there is
\(\lambda(\varepsilon,n)<\infty\) such that a complete nonnegatively curved
\((\mathcal M^n,g)\) with basepoint \(O\) contains only finitely many disjoint
balls \(B(p,r)\) satisfying
\[
 \operatorname{sect}(g)\ge\varepsilon r^{-2}\ \text{on }B(p,r),
 \qquad d(p,O)\ge\lambda r.
\]
\end{theorem}

\section{Space-time point picking with restrictions}

For a Ricci flow \((\mathcal N^n,h(t))\), \(t\in[t_0,0]\), with
\(\Ric\ge0\), and constants \(B,C>0\), put
\[
 \mathcal N_{B,C}=\{(y,t):t>t_0, R(y,t)>C+B(t-t_0)^{-1}\}.
\]

\begin{lemma}[Existence of a large-curvature point with local control]
\label{lem:chow-iii-spacetime-point-picking}
\source{chow-et-al-ricci-flow-part-iii:page-0078,chow-et-al-ricci-flow-part-iii:page-0079,chow-et-al-ricci-flow-part-iii:page-0080}
\dcref{ch18:18.24}
\group{Space-Time Point Picking}

Suppose \(B_{h(0)}(p,1)\Subset\mathcal N\), and let
\((y_1,t_1)\in\mathcal N_{B,C}\) satisfy
\(y_1\in B_{h(t_1)}(p,1/4)\).  For \(\sigma\in(0,1)\), set
\[
 A_1(\sigma)=\bigl(C+B(t_1-t_0)^{-1}\bigr)
 \min\left\{\frac\sigma2(t_1-t_0),\frac1{2304}\right\}.
\]
If \(A_1(\sigma)>1\), then for every \(A\in(1,A_1(\sigma)]\) there is
\((y_*,t_*)\in\mathcal N_{B,C}\), with
\(t_*-t_0\ge(1-\sigma)(t_1-t_0)\) and
\(y_*\in B_{h(t_*)}(p,1/3)\), such that \(R(y,t)\le2R(y_*,t_*)\) whenever
\[
 t\in(t_*-A/R(y_*,t_*),t_*],
 \quad d_{h(t)}(y,p)\le d_{h(t_*)}(y_*,p)+\sqrt{A/R(y_*,t_*)},
\]
and \((y,t)\in\mathcal N_{B,C}\).
\end{lemma}
\begin{proof}
\sketch Starting at \((y_1,t_1)\), replace \((y_i,t_i)\) by an admissible point of
more than twice its scalar curvature whenever the conclusion fails.  The
curvatures then grow geometrically.  Consequently the accumulated spatial
displacement is less than
\(4A^{1/2}(C+B(t_1-t_0)^{-1})^{-1/2}\le1/12\), while the accumulated backward
time is at most \(2A(C+B(t_1-t_0)^{-1})^{-1}\le\sigma(t_1-t_0)\).  Thus every
chosen point stays in the compact cylinder
\(B_{h(0)}(p,1/3)\times[t_0,t_1]\), where curvature is bounded.  Geometric
growth therefore forces the process to terminate, and its final point has all
the asserted properties.
\end{proof}

\begin{proposition}[Curvature control in a parabolic cylinder]
\label{prop:chow-iii-parabolic-cylinder-curvature-control}
\source{chow-et-al-ricci-flow-part-iii:page-0080,chow-et-al-ricci-flow-part-iii:page-0081,chow-et-al-ricci-flow-part-iii:page-0082,chow-et-al-ricci-flow-part-iii:page-0083}
\uses{lem:chow-iii-spacetime-point-picking, thm:chow-iii-changing-distances}
\dcref{ch18:18.25}
\group{Space-Time Point Picking}


Let \((y_*,t_*)\) be as in
Lemma~\ref{lem:chow-iii-spacetime-point-picking}, and set
\[
 D=\min\left\{\frac{A^{1/2}}{100(n-1)},
 \frac{C(t_*-t_0)+B}{2}\right\}.
\]
Then \(R(y,t)\le2R(y_*,t_*)\) on
\[
 B_{h(t_*)}\left(y_*,\frac1{10}A^{1/2}R(y_*,t_*)^{-1/2}\right)
 \times[t_*-D/R(y_*,t_*),t_*].
\]
\end{proposition}
\begin{proof}
\sketch If a violating point existed, it would lie in \(\mathcal N_{B,C}\).  Follow
the distance from that point to \(p\) backward until it first reaches half the
allowed spatial margin.  On this interval all endpoint balls used by
Theorem~\ref{thm:chow-iii-changing-distances} have scalar curvature at most
\(2R(y_*,t_*)\), either by the definition of \(\mathcal N_{B,C}\) or by
Lemma~\ref{lem:chow-iii-spacetime-point-picking}.  Integrating the changing
distance estimate with \(r_0=R(y_*,t_*)^{-1/2}/10\) shows that the distance can
increase by less than the available margin.  This contradicts the first-time
choice, so no violating point exists.
\end{proof}

\section{Necks in manifolds with positive sectional curvature}

Put \(c_n=\sqrt{(n-1)(n-2)}\).

\begin{definition}[Embedded epsilon-neck]
\label{def:chow-iii-embedded-epsilon-neck}
\source{chow-et-al-ricci-flow-part-iii:page-0083,chow-et-al-ricci-flow-part-iii:page-0084}
\dcref{ch18:18.26}
\group{Necks}

Let \((\mathcal N^n,h)\) be complete, \(R(p)>0\), and
\(r=c_nR(p)^{-1/2}\).  The ball \(\mathfrak N=B(p,\varepsilon^{-1}r)\) is an
embedded \(\varepsilon\)-neck if there is an embedding
\(\varphi:\mathfrak N\to S^{n-1}\times\mathbb R\), with
\(\varphi(p)\in S^{n-1}\times\{0\}\), such that
\[
 \left|\varphi^*(g_{S^{n-1}}+\dd u^2)-r^{-2}h\right|_
 {C^{\lceil\varepsilon^{-1}\rceil+1}(r^{-2}h)}<\varepsilon.
\]
The inverse image of \(S^{n-1}\times\{0\}\) is its center sphere.
\end{definition}

\begin{remark}[Comparison with fixed-quality necks]
\label{rem:chow-iii-neck-quality-comparison}
\source{chow-et-al-ricci-flow-part-iii:page-0084}
\uses{def:chow-iii-embedded-epsilon-neck}
\dcref{ch18:18.27}
\group{Necks}


Every embedded \(\varepsilon\)-neck with
\(\varepsilon\le\varepsilon(n)\) is also an embedded
\(\varepsilon(n)\)-neck.  Given fixed \((\varepsilon_0,k_0,L_0)\), sufficiently
small \(\varepsilon\) also guarantees the presence of an embedded neck of
that fixed quality and length.
\end{remark}

\begin{lemma}[Center spheres of sufficiently good necks bound smooth balls]
\label{lem:chow-iii-neck-center-bounds-ball}
\source{chow-et-al-ricci-flow-part-iii:page-0084}
\uses{def:chow-iii-embedded-epsilon-neck}
\dcref{ch18:18.28}
\group{Necks}


There is \(\varepsilon(n)>0\) such that, if \((\mathcal N^n,h)\) is complete,
noncompact, and has positive sectional curvature, then the center sphere of
every embedded \(\varepsilon(n)\)-neck bounds a smooth \(n\)-ball in
\(\mathcal N\).
\end{lemma}

\begin{lemma}[Diameter of a Busemann level through a neck]
\label{lem:chow-iii-busemann-level-diameter-neck}
\source{chow-et-al-ricci-flow-part-iii:page-0085,chow-et-al-ricci-flow-part-iii:page-0086,chow-et-al-ricci-flow-part-iii:page-0087,chow-et-al-ricci-flow-part-iii:page-0088}
\uses{lem:chow-iii-neck-center-bounds-ball}
\dcref{ch18:18.30}
\group{Necks and Busemann Functions}


For \(n\ge3\), there is \(\varepsilon(n)>0\) with the following property.
Let \((\mathcal N^n,h)\) be complete, noncompact, and positively curved; let
\(\gamma\) be a ray; and let \(q_0\) lie on the center sphere of an embedded
\(\varepsilon(n)\)-neck.  If \(b_0=b_\gamma(q_0)\), then
\[
 \frac{11}{12}\pi c_nR(q_0)^{-1/2}
 \le\operatorname{diam}b_\gamma^{-1}(b_0)
 \le11\pi c_nR(q_0)^{-1/2}.
\]
\end{lemma}
\begin{proof}
\sketch Lemma~\ref{lem:chow-iii-neck-center-bounds-ball} separates the neck into an
inner compact side and an outer side.  Cylinder closeness and Toponogov
comparison imply that \(b_\gamma<b_0\) beyond a fixed distance on the inner
side and \(b_\gamma>b_0\) beyond that distance on the outer side.  Thus the
level set lies in the cylindrical slab \(S^{n-1}\times[-5\pi,5\pi]\), whose
diameter, after allowing the prescribed neck error, is at most the displayed
upper bound.  For each \(z\in S^{n-1}\), continuity gives a point
\((z,u(z))\) in that slab on the level set.  Taking \(z\) antipodal to the
angular coordinate of \(q_0\), cylinder closeness gives the displayed lower
bound.
\end{proof}

\begin{lemma}[A sufficiently good neck forces a Busemann minimum]
\label{lem:chow-iii-busemann-attains-minimum}
\source{chow-et-al-ricci-flow-part-iii:page-0088}
\uses{lem:chow-iii-busemann-level-diameter-neck}
\dcref{ch18:18.31}
\group{Necks and Busemann Functions}


There is \(\varepsilon(n)>0\) such that, if a complete noncompact positively
curved \((\mathcal N^n,h)\) contains an embedded \(\varepsilon(n)\)-neck,
then the Busemann function of every ray is bounded below and attains its
minimum.
\end{lemma}
\begin{proof}
\sketch Lemma~\ref{lem:chow-iii-busemann-level-diameter-neck} gives a compact
closure for the union of the inner component and the neck, while the Busemann
function on its complement exceeds its value at a center-sphere point.  Its
global infimum is therefore achieved on that compact set.
\end{proof}

\begin{remark}[The neck hypothesis is essential]
\label{rem:chow-iii-busemann-minimum-neck-essential}
\source{chow-et-al-ricci-flow-part-iii:page-0088}
\dcref{ch18:18.32}
\group{Necks and Busemann Functions}

Without a sufficiently good neck, a Busemann function need not be bounded
below.
\end{remark}

\begin{proposition}[Farther necks have almost larger radii]
\label{prop:chow-iii-farther-necks-larger-radii}
\source{chow-et-al-ricci-flow-part-iii:page-0089}
\uses{lem:chow-iii-busemann-level-diameter-neck, lem:chow-iii-busemann-attains-minimum}
\dcref{ch18:18.33}
\group{Necks and Busemann Functions}


There is \(\varepsilon(n)>0\) such that the following holds.  Let
\((\mathcal N^n,h)\) be complete, noncompact, and positively curved, let
\(\gamma\) be a ray, and let \(\mathfrak N_1,\mathfrak N_2\) be disjoint
embedded \(\varepsilon(n)\)-necks with center points \(y_1,y_2\).  If
\(d(p_*,y_1)\le d(p_*,y_2)\), where \(p_*\) minimizes \(b_\gamma\), then
\[
 R(y_2)\le144R(y_1).
\]
\end{proposition}
\begin{proof}
\sketch The neck separation inequalities give \(b_\gamma(y_1)<b_\gamma(y_2)\).
Sharafutdinov monotonicity for the convex Busemann function then gives
\(\operatorname{diam}b_\gamma^{-1}(b_\gamma(y_1))\le
\operatorname{diam}b_\gamma^{-1}(b_\gamma(y_2))\).  Applying the lower and
upper diameter bounds of Lemma~\ref{lem:chow-iii-busemann-level-diameter-neck}
at \(y_1\) and \(y_2\), respectively, gives
\((11/12)\pi c_nR(y_1)^{-1/2}\le11\pi c_nR(y_2)^{-1/2}\), which is equivalent
to the assertion.
\end{proof}

\section{Localized no local collapsing}

\begin{definition}[Weakly kappa-noncollapsed]
\label{def:chow-iii-weakly-kappa-noncollapsed}
\source{chow-et-al-ricci-flow-part-iii:page-0090}
\dcref{ch18:18.34}
\group{Localized Noncollapsing}

A complete bounded-curvature Ricci flow \(g(t)\), \(t\in[0,T)\), is weakly
\(\kappa\)-noncollapsed at \((x_0,t_0)\) below scale \(\rho\) if, whenever
\(0<r\le\rho\) and
\[
 |\Rm|\le r^{-2}\quad\text{on }B_{g(t_0)}(x_0,r)\times[t_0-r^2,t_0],
\]
one has \(\Vol_{g(t_0)}B_{g(t_0)}(x_0,r)\ge\kappa r^n\).
\end{definition}

\begin{remark}[Largest curvature-controlled scale]
\label{rem:chow-iii-largest-curvature-controlled-scale}
\source{chow-et-al-ricci-flow-part-iii:page-0090}
\uses{def:chow-iii-weakly-kappa-noncollapsed}
\dcref{ch18:18.35}
\group{Localized Noncollapsing}


The function
\[
 \Psi(r)=r^2\max_{B_{g(t_0)}(x_0,r)\times[t_0-r^2,t_0]}|\Rm|
\]
is nondecreasing and tends to zero with \(r\).  If \(\bar r\) is the largest
radius on which \(\Psi\le1\), weak \(\kappa\)-noncollapsing below \(\rho\) is
equivalent to the volume lower bound for all
\(0<r\le\min\{\rho,\bar r\}\).
\end{remark}

\begin{theorem}[Localized no local collapsing]
\label{thm:chow-iii-localized-no-local-collapsing}
\source{chow-et-al-ricci-flow-part-iii:page-0090,chow-et-al-ricci-flow-part-iii:page-0091,chow-et-al-ricci-flow-part-iii:page-0092,chow-et-al-ricci-flow-part-iii:page-0093}
\uses{def:chow-iii-weakly-kappa-noncollapsed, lem:chow-iii-localized-reduced-distance-control}
\dcref{ch18:18.36}
\group{Localized Noncollapsing}


For every \(n\ge2\) and \(A>0\), there is \(\kappa(n,A)>0\) such that the
following holds.  Let \(g(t)\), \(t\in[0,r_0^2]\), be a complete
bounded-curvature Ricci flow with
\[
 |\Rm|\le r_0^{-2}\quad\text{on }B_{g(0)}(x_0,r_0)\times[0,r_0^2],
 \qquad \Vol_{g(0)}B_{g(0)}(x_0,r_0)\ge A^{-1}r_0^n.
\]
Then at time \(r_0^2\), the flow is weakly \(\kappa(n,A)\)-noncollapsed below
scale \(r_0\) at every point of \(B_{g(r_0^2)}(x_0,Ar_0)\).
\end{theorem}
\begin{proof}
\sketch Scale to \(r_0=1\), reverse time, and base reduced distance and reduced volume
at a candidate point \(x\in B_{g(1)}(x_0,A)\).  If a radius \(r\le1\) has
curvature at most \(r^{-2}\), write
\(\kappa=\Vol_{g(1)}B_{g(1)}(x,r)/r^n\).  The standard reduced-volume estimate
at the collapsed scale gives
\(\widetilde V(\kappa^{1/n}r^2)\le c_2(n,\kappa)\), where
\(c_2(n,\kappa)\to0\) as \(\kappa\downarrow0\).

Lemma~\ref{lem:chow-iii-localized-reduced-distance-control} supplies a point
at backward time \(1/2\) with uniformly bounded reduced distance.  Concatenate
a minimizing reduced geodesic to that point with a controlled ordinary
geodesic inside the unit ball.  The curvature and metric-comparison estimates
give \(\ell(q,1)\le c_5(n,A)\) throughout \(B_{g(1)}(x_0,1)\).  Hence
\[
 \widetilde V(1)\ge(4\pi)^{-n/2}e^{-c_5(n,A)}A^{-1}=:c_3(n,A)>0.
\]
Reduced-volume monotonicity yields
\(c_3(n,A)\le c_2(n,\kappa)\), which prevents \(\kappa\) from tending to zero
and furnishes the required \(\kappa(n,A)\).
\end{proof}

\begin{remark}[The theorem at the maximal initial scale]
\label{rem:chow-iii-localized-noncollapsing-maximal-scale}
\source{chow-et-al-ricci-flow-part-iii:page-0090}
\uses{thm:chow-iii-localized-no-local-collapsing}
\dcref{ch18:18.37}
\group{Localized Noncollapsing}


For a nonflat flow on \([0,T]\), let \(\bar r_0\le\sqrt T\) be the largest
radius for which the curvature bound in
Theorem~\ref{thm:chow-iii-localized-no-local-collapsing} holds, and set
\(A(r_0)=r_0^n/\Vol_{g(0)}B_{g(0)}(x_0,r_0)\).  For every
\(r_0\le\bar r_0\), the theorem gives weak noncollapsing below scale \(r_0\)
on \(B_{g(r_0^2)}(x_0,A(r_0)r_0)\).
\end{remark}

\begin{lemma}[Localized reduced-distance control]
\label{lem:chow-iii-localized-reduced-distance-control}
\source{chow-et-al-ricci-flow-part-iii:page-0092,chow-et-al-ricci-flow-part-iii:page-0093,chow-et-al-ricci-flow-part-iii:page-0094,chow-et-al-ricci-flow-part-iii:page-0095,chow-et-al-ricci-flow-part-iii:page-0096,chow-et-al-ricci-flow-part-iii:page-0097}
\uses{thm:chow-iii-changing-distances}
\dcref{ch18:18.38}
\group{Localized Noncollapsing}


Let \(g(t)\), \(t\in[0,1]\), be a complete bounded-curvature Ricci flow, and
suppose
\(\lvert\Rm\rvert\le1\) on \(B_{g(0)}(x_0,1)\times[0,1]\).  Then
\begin{enumerate}
\item \(|W|_{g(t)}\le e^{n-1}|W|_{g(0)}\) for
\(z\in B_{g(0)}(x_0,1)\), \(W\in T_z\mathcal M\), and \(t\in[0,1]\);
\item \(B_{g(t)}(x_0,e^{1-n})\subset B_{g(0)}(x_0,1)\);
\item in the setting of Theorem~\ref{thm:chow-iii-localized-no-local-collapsing}
with \(r_0=1\), there is \(c_4(n,A)<\infty\) such that
\[
 \min_{y\in\overline B_{g(1/2)}(x_0,e^{1-n})}\ell(y,1/2)\le c_4(n,A).
\]
\end{enumerate}
\end{lemma}
\begin{proof}
\sketch The curvature bound gives \(|\Ric|\le n-1\); integrating
\(|\partial_tg(W,W)|\le2(n-1)g(W,W)\) proves the first assertion, and applying
it along minimizing geodesics proves the second.

For the third, choose a nondecreasing cutoff \(\phi\) which is one below
\(e^{1-n}/2\), diverges at \(e^{1-n}\), and satisfies
\[
 2(\phi')^2/\phi-\phi''\ge(2A+ne^n)\phi'-C_1(n,A)\phi.
\]
For backward time \(\tau=1-t\), put
\[
 h(w,t)=\phi(d_{g(t)}(x_0,w)-(2t-1)A)(\overline L(w,1-t)+2n+1).
\]
The scalar-curvature lower bound gives \(\overline L\ge-2n\) for
\(\tau\le1/2\).  At a spatial minimum of \(h\), the reduced-distance
inequality \(\overline L_t+\Delta\overline L\le2n\), the cutoff inequality,
and Theorem~\ref{thm:chow-iii-changing-distances} imply
\[
 (\partial_t-\Delta)h\ge-(2n+C_1(n,A))h.
\]
The minimum principle, integrated from \(t=1\) to \(t=1/2\), bounds the
minimum of \(h(\cdot,1/2)\).  Since \(\phi\) diverges at the boundary of the
ball and is at least one, that minimum occurs inside
\(B_{g(1/2)}(x_0,e^{1-n})\) and yields the asserted bound for \(\ell\).
\end{proof}

\begin{remark}[Nonsmoothness of the localized barrier]
\label{rem:chow-iii-localized-barrier-nonsmoothness}
\source{chow-et-al-ricci-flow-part-iii:page-0097}
\dcref{ch18:18.39}
\group{Localized Noncollapsing}

The nonsmooth distance terms in the proof of
Lemma~\ref{lem:chow-iii-localized-reduced-distance-control} are interpreted
using upper support functions and the corresponding barrier maximum
principle.
\end{remark}

\begin{example}[Problem: dependence on the minimizing geodesic]
\label{ex:chow-iii-distance-derivative-any-minimizer}
\dcref{ch18:18.3}
\group{Changing Distances}
\source{chow-et-al-ricci-flow-part-iii:page-0063}
Determine whether, under Ricci flow, the formula in
Corollary~\ref{cor:chow-iii-forward-distance-derivative-ricci-flow} equals
\(-\int_\gamma\Ric(\gamma',\gamma')\,\dd s\) for every
\(\gamma\in\mathcal Z(t_0)\).
\end{example}

\begin{example}[Problem: general point picking]
\label{ex:chow-iii-general-point-picking}
\dcref{ch18:18.40}
\group{Notes and Problems}
\source{chow-et-al-ricci-flow-part-iii:page-0097}
Formulate general point-picking results for a metric space \(S\), an objective
function \(f:S\to\mathbb R\), and a reference function
\(\phi:S\to\mathbb R\), encompassing the spatial and space-time methods of
this chapter.
\end{example}

\begin{example}[Conjecture: smooth Schoenflies in dimension four]
\label{ex:chow-iii-smooth-schoenflies-four}
\dcref{ch18:18.29}
\group{Necks}
\source{chow-et-al-ricci-flow-part-iii:page-0084}
Every smoothly embedded three-sphere in \(\mathbb R^4\) bounds a smooth
four-ball.
\end{example}

\begin{example}[Exercise: trivial pointed limits after rescaling]
\label{ex:chow-iii-trivial-limits-rescaling}
\dcref{ch18:18.41}
\group{Notes and Problems}
\source{chow-et-al-ricci-flow-part-iii:page-0097}
For every sequence of pointed Riemannian manifolds
\((\mathcal M_i^n,g_i,O_i)\), there are \(K_i\to\infty\) such that
\((\mathcal M_i^n,K_ig_i,O_i)\) converges to pointed Euclidean space.
\end{example}
\begin{proof}
Choose a normal-coordinate radius \(\rho_i>0\) at \(O_i\) on which the metric
and its first \(i\) derivatives are within \(1/i\) of the Euclidean metric
after identifying \(T_{O_i}\mathcal M_i\) with \(\mathbb R^n\).  Choose
\(K_i\ge i\) so large that every fixed-radius ball in \(K_ig_i\) lies inside
that normal-coordinate ball and the rescaled curvature derivatives through
order \(i\) are at most \(1/i\).  A diagonal argument then gives smooth
pointed convergence to \((\mathbb R^n,g_{\mathrm{Euc}},0)\).
\end{proof}
